\documentclass{article}

% Recommended, but optional, packages for figures and better typesetting:
\usepackage{microtype}
\usepackage{graphicx}
\usepackage{subcaption}
\usepackage{booktabs} % for professional tables
\usepackage{hyperref}
\newcommand{\theHalgorithm}{\arabic{algorithm}}
\usepackage{icml2026}

\usepackage{amsmath}
\usepackage{amssymb}
\usepackage{mathtools}
\usepackage{amsthm}
\usepackage[capitalize,noabbrev]{cleveref}

% For pseudocode
\usepackage{algorithm}
\usepackage{algorithmic}

%%%%%%%%%%%%%%%%%%%%%%%%%%%%%%%%
% THEOREMS
%%%%%%%%%%%%%%%%%%%%%%%%%%%%%%%%
\theoremstyle{plain}
\newtheorem{theorem}{Theorem}[section]
\newtheorem{proposition}[theorem]{Proposition}
\newtheorem{lemma}[theorem]{Lemma}
\newtheorem{corollary}[theorem]{Corollary}
\theoremstyle{definition}
\newtheorem{definition}[theorem]{Definition}
\newtheorem{assumption}[theorem]{Assumption}
\theoremstyle{remark}
\newtheorem{remark}[theorem]{Remark}

\icmltitlerunning{Robust Open-World Test-Time Model Merging}

\begin{document}

\twocolumn[
  \icmltitle{Robust Open-World Test-Time Model Merging via \\
    Dynamic Fisher Preconditioning and Dynamic Feature Centering}

  \begin{icmlauthorlist}
    \icmlauthor{Anonymous Author}{equal,dept}
  \end{icmlauthorlist}

  \icmlaffiliation{dept}{Department of Computer Science, Anonymous University, Location, Country}
  \icmlcorrespondingauthor{Anonymous Author}{author@univ.edu}

  \icmlkeywords{Model Merging, Test-Time Adaptation, Open-World, Fisher Preconditioning, Contrast Invariance}

  \vskip 0.3in
]

\printAffiliationsAndNotice{}

\begin{abstract}
  Test-Time Model Merging (TTMM) is an emerging paradigm to dynamically combine specialized expert networks in weight space on-the-fly to handle non-stationary, unlabeled test streams. Existing state-of-the-art open-world TTMM frameworks (such as DR-Fisher) leverage class prototype cohesion for novelty detection and entropy-based routing to handle out-of-distribution data. However, we identify two critical limitations in prior frameworks under realistic environmental corruptions like contrast reduction and noise: (1) they rely on offline static feature centering, which we show is extremely vulnerable to global feature scaling, causing catastrophic novelty detection collapse (100\% False Positive Rate and 10\% Novelty Detection Rate under contrast); and (2) they employ a frozen test-time preconditioning matrix that cannot adapt to shifting test manifolds. To resolve these challenges, we propose \textbf{Robust Open-World Test-Time Model Merging with Dynamic Fisher and Dynamic Feature Centering (D-TT-Fisher-DFC)}. First, we introduce \textbf{Dynamic Feature Centering (DFC)}, which dynamically computes and subtracts the mean of the features on the active test batch. We mathematically prove that DFC makes the prototype-cohesion metric invariant to global contrast scaling. Second, we maintain a running, exponentially-smoothed estimate of the test-time diagonal Fisher sensitivities to adaptively scale learning rates layer-wise. Our evaluations on non-stationary, corrupted vision streams demonstrate that our framework achieves perfect novelty detection (100\% NDR, 0\% FPR under clean/noise, and 100\% NDR, 10\% FPR under contrast) while maintaining near-perfect classification accuracy, outperforming the state-of-the-art by up to +90\% absolute under corruption.
\end{abstract}

\section{Introduction}
\label{sec:intro}
Deep neural networks are frequently deployed in non-stationary environments where the test data distribution shifts over time due to changes in weather, illumination, sensor noise, or task transitions \cite{liang2023comprehensive, niu2023stable, hendrycks2019benchmarking}. In real-world applications such as autonomous driving, medical imaging, and mobile robotics, systems must operate reliably under unpredictable environmental conditions without the luxury of costly re-training or source calibration data. Traditionally, Fully Test-Time Adaptation (TTA) fine-tunes a single pre-trained model on unlabeled streaming data on-the-fly \cite{wang2021tent}. While elegant, standard TTA often suffers from catastrophic forgetting and error accumulation when exposed to diverse multi-task distributions or non-stationary streams where the target domain shifts repeatedly.

To address this multi-task challenge, \emph{Test-Time Model Merging (TTMM)} has emerged as a powerful, training-free alternative \cite{zhao2024dynamic, yang2024adamerging}. Instead of modifying network weights directly through expensive gradient steps across the entire backbone, TTMM dynamically blends multiple specialized task-specific experts $\{\theta_1, \ldots, \theta_K\}$ in parameter space on-the-fly. By leveraging the convex combination of task-specific parameter vectors (task arithmetic) \cite{ilharco2023editing, wortsman202ModelSoups}, TTMM can swiftly navigate complex multi-task streams without suffering from catastrophic forgetting. Because each expert network remains isolated and is only combined in the parameter weight space at test-time, the model retains sharp knowledge boundaries.

Despite its promise, real-world deployment requires operating in an \emph{open-world} setting. In this scenario, the incoming test stream is not restricted to known domains but contains both known domains (for which specialized experts are available) and completely novel, unseen domains \cite{proto_ttmm}. Existing open-world TTMM methods (e.g., PROTO-TTMM \cite{proto_ttmm}, IGGS-OW \cite{iggs_ow}, and DR-Fisher \cite{dr_fisher}) utilize isotropic feature centering and class prototype cohesion to detect whether an active batch belongs to a known task or a novel task. While effective under pristine, clean conditions, we find that this mechanism fails catastrophically under realistic environmental corruptions such as noise and contrast reduction. These corruptions are physically common (e.g., fog, rain, lens dirt, or low-light conditions) and severely alter the feature scale.

Specifically, we identify two major scientific gaps in existing state-of-the-art open-world TTMM pipelines:
\begin{enumerate}
  \item \textbf{Vulnerability of Static Centering to Scale Shifts:} Prior work precomputes dataset mean features (e.g., $\mu_{\text{static}}$) on clean calibration data and uses them to center online test features. Under contrast reduction or illumination changes, test features undergo global scaling (e.g., scaled by $c = 0.3$). We mathematically prove that subtracting a static clean mean from scaled test features introduces a severe directional bias towards $-\mu_{\text{static}}$. This directional bias completely destroys the semantic class structures and collapses prototype cohesion, leading to a 100\% False Positive Rate (falsely flagging all known digits as novel) and a 10\% Novelty Detection Rate (failing to detect clothing as novel).
  \item \textbf{Frozen Preconditioning Sensitivities:} Existing preconditioned methods like DR-Fisher \cite{dr_fisher} or FP-OW \cite{fp_ow} estimate parameter sensitivities (TT-Fisher) once on an initial calibration window and freeze them. Under highly non-stationary streams where corruptions and domains shift mid-stream, frozen sensitivities become misaligned with the active data manifold.
\end{enumerate}

To overcome these fundamental limits, we propose \textbf{Robust Open-World Test-Time Model Merging via Dynamic Fisher Preconditioning and Dynamic Feature Centering (D-TT-Fisher-DFC)}. Our framework makes two core contributions:
\begin{itemize}
  \item \textbf{Dynamic Feature Centering (DFC):} Instead of relying on precomputed offline statistics, DFC centers test features using the mean of the active test batch on-the-fly. We mathematically prove that DFC renders the prototype-cohesion metric completely invariant to any global scale factor $c > 0$ (such as contrast scaling).
  \item \textbf{Dynamic Test-Time Fisher (D-TT-Fisher):} We maintain a running exponential moving average (EMA) of the diagonal Fisher sensitivities over streaming batches, providing true online preconditioning that dynamically adjusts layer-wise learning rates to shifting domains and corruptions.
\end{itemize}

Empirically, D-TT-Fisher-DFC demonstrates outstanding robustness across all vision streams (MNIST, KMNIST, FashionMNIST) under clean, noise, and contrast corruptions. It completely eliminates false positives (0\% FPR under clean and noise, 10\% FPR under contrast) and achieves near-perfect novelty detection (100\% NDR under contrast), improving upon the state-of-the-art by up to \textbf{+90\% absolute} in novelty detection quality, while preserving perfect multi-task classification accuracy.

\section{Related Work}
\label{sec:related}

\subsection{Model Merging and Task Arithmetic}
Model merging combines multiple task-specific expert neural networks into a single multi-task network without additional training or access to raw training data \cite{wortsman202ModelSoups, matena2022merging}. Initial methods relied on simple weight averaging or Stochastic Weight Averaging (SWA) \cite{izmailov2018averaging}. To resolve weight-space permutations, Git Re-basin \cite{ainsworth2023git} aligns weights before merging. Advanced merging methods like RegMean \cite{jin2023regmean}, Ties-Merging \cite{yadav2023ties}, DARE \cite{yu2023dare}, and ZipIt! \cite{stoica2024zipit} address parameter interference and representation mismatch during weight fusion. 

Recently, editing models using task arithmetic has emerged as a prominent paradigm \cite{ilharco2023editing}, and recent 2024-2025 works have introduced advanced techniques to reduce task interference, such as Task Singular Vectors \cite{gargiulo2025task}, guiding data-free model merging via task vectors \cite{cheng2025whoever}, sparse fine-tuning \cite{iurada2025efficient}, uncertainty-based gradient matching \cite{daheim2024model_gradient}, knowledge editing \cite{fu2025model}, and localizing task information \cite{wang2024localizing}. Furthermore, model merging is finding applications in federated learning \cite{morafah2024towards} and has been surveyed extensively \cite{yang2024model_survey}. However, almost all of these methods are offline and static, assuming a fixed combination of tasks, making them unable to adapt to dynamic, non-stationary test-time streams.

\subsection{Test-Time Adaptation (TTA)}
Test-Time Adaptation (TTA) adapts a pre-trained model to distribution shifts using unlabeled streaming data \cite{wang2021tent, niu2022efficient}. Fully test-time adaptation methods like Tent \cite{wang2021tent} minimize prediction entropy to update normalization layers. To handle non-stationary, continuously shifting environments, Continual TTA (CoTTA) \cite{wang2022continual_cotta} and EATA \cite{niu2022efficient_eata} introduce anti-forgetting and selection mechanisms. Test-Time Training (TTT) \cite{sun2020test_ttt} utilizes self-supervised auxiliary tasks. Recent works focus on robust TTA under noisy data streams (SoTTA \cite{gong2023sotta}), malicious test samples (MedBN \cite{Park_2024_CVPR}), and dynamic scenarios (Robust TTA \cite{yuan2023robust_dynamic}). Comprehensive overviews of the TTA landscape are provided in recent surveys \cite{liang2024comprehensive, xiao2024beyond, wang2023search}. However, standard TTA methods modify the entire network backbone, making them prone to representation collapse and catastrophic forgetting under multi-task shifts, whereas model merging naturally preserves task boundaries.

\subsection{Open-World Test-Time Model Merging}
To combine the multi-task capacity of model merging with the dynamic adaptation of TTA, recent works propose Test-Time Model Merging (TTMM) \cite{zhao2024dynamic, yang2024adamerging}. To break the closed-world assumption where only known tasks appear, PROTO-TTMM \cite{proto_ttmm} and IGGS-OW \cite{iggs_ow} introduce class prototype cohesion to detect novel tasks. DR-Fisher \cite{dr_fisher} resolves Batch Normalization (BN) activation mismatch during merging and uses Entropy-Based Expert Routing (EBER) to bypass the feedback loop trap. However, all these frameworks remain extremely vulnerable to environmental corruptions like noise and contrast shifts, which we analyze and resolve in this paper.

\section{Preliminaries and Problem Formulation}
\label{sec:prelims}
We consider an open-world Test-Time Model Merging (TTMM) setting \cite{proto_ttmm}. We are given $K$ specialized expert models $\{\theta_1, \ldots, \theta_K\}$ sharing a common pre-trained backbone and a base initialization $\theta_{\text{base}}$. Each expert $\theta_k$ is fine-tuned on a distinct known source domain $\mathcal{D}_k$. The task vectors are defined as $v_k = \theta_k - \theta_{\text{base}}$. At test time, we receive a continuous, non-stationary stream of unlabeled batches $B_1, B_2, \ldots, B_T$, where each batch $B_t = \{x_i\}_{i=1}^B$ can be from a known source domain $\mathcal{D}_k$ or a completely novel, unseen domain $\mathcal{D}_{\text{novel}}$. The objective is to dynamically solve for layer-wise merging coefficients $\Lambda_t^{(w)} \in \mathbb{R}^K$ for each parameter tensor $w$ to assemble a merged model:
\begin{equation}
  \theta_{\text{merged}} = \theta_{\text{base}} + \sum_{k=1}^K \lambda_{t, k}^{(w)} v_k
\end{equation}
that maximizes classification accuracy on the active batch, while robustly identifying novel domains and routing known domains.

\textbf{Unified Static Space Precomputation:} To avoid representational space mismatch between different experts, we use a uniform static anchor model $\theta_{\text{static}} = \theta_{\text{base}} + \frac{1}{K}\sum_k v_k$. The offline class prototypes $P_k = \{\pi_{k,c}\}_{c=0}^{C-1}$ represent the mean centered features of class $c$ in domain $\mathcal{D}_k$ extracted using $\theta_{\text{static}}$.

\textbf{Static Isotropic Centering:} To compute cohesion, prior works center the test features $f(x)$ using the precomputed offline clean mean $\mu_k = \mathbb{E}_{x \sim \mathcal{D}_k}[f(x; \theta_{\text{static}})]$:
\begin{equation}
  z_i = f(x_i; \theta_{\text{static}}) - \mu_k
\end{equation}
The cohesion score of a batch $B_t$ to expert $k$ is defined as the average maximum cosine similarity:
\begin{equation}
  C_k(B_t) = \frac{1}{B} \sum_{i=1}^B \max_{c \in \{0, \dots, C-1\}} \text{sim}(z_i, \pi_{k,c})
\end{equation}
where $\text{sim}(u, v) = \frac{u^T v}{\|u\|_2 \|v\|_2}$ is the cosine similarity. If $\max_k C_k(B_t) < \tau_N$, where $\tau_N$ is a pre-defined novelty threshold, the batch is flagged as novel ($\text{is\_novel} = \text{True}$). Otherwise, the batch is routed to the best expert $k^* = \arg\max_k C_k(B_t)$.

\textbf{Entropy-Based Expert Routing (EBER):} Rather than optimizing coefficients from scratch, EBER routes the batch to the expert that minimizes predictive entropy:
\begin{equation}
  k^* = \arg\min_k \frac{1}{B} \sum_{i=1}^B H(p(y \mid x_i; \theta_k))
\end{equation}
where $H(p) = -\sum_c p_c \log p_c$. EBER sets the routing prior weight for $k^*$ to $0.99$ and all other experts to $0.005$, creating a highly accurate routing prior that bypasses the feedback loop trap of joint optimization.

\section{Proposed Method: D-TT-Fisher-DFC}
\label{sec:method}
In this section, we present our two core innovations: \textbf{Dynamic Feature Centering (DFC)} to make open-world novelty routing robust to global scale corruptions, and \textbf{Dynamic Test-Time Fisher (D-TT-Fisher)} to provide adaptive online preconditioning.

\subsection{Dynamic Feature Centering (DFC)}
Under environmental corruptions such as contrast reduction, the pixel intensities are scaled down. Due to the linear nature of early neural network layers, the extracted feature vectors also undergo a global scaling. Let $f(x) \approx c \cdot f_{\text{clean}}(x)$, where $c \in (0, 1]$ represents the contrast scale factor (e.g., $c = 0.3$).

Under static centering, the centered feature is computed as:
\begin{align}
  z_i^{\text{static}} &= f(x_i) - \mu_{\text{clean}} \nonumber \\
  &\approx c \cdot f_{\text{clean}}(x_i) - \mu_{\text{clean}} \nonumber \\
  &= c(f_{\text{clean}}(x_i) - \mu_{\text{clean}}) + (c - 1)\mu_{\text{clean}}
\end{align}
Since $c \in (0, 1)$, the second term $(c - 1)\mu_{\text{clean}}$ is non-zero and points in the direction of $-\mu_{\text{clean}}$. Under severe contrast reduction (e.g., $c = 0.3$), the bias term $-0.7\mu_{\text{clean}}$ completely dominates the centered feature, causing all $z_i^{\text{static}}$ to align heavily with $-\mu_{\text{clean}}$ regardless of their actual class. This destroys class separation, rendering the cosine similarity with class prototypes meaningless and causing novelty detection to collapse.

To resolve this, we propose \textbf{Dynamic Feature Centering (DFC)}. DFC computes and subtracts the mean of the features on the active test batch $B_t$ on-the-fly:
\begin{equation}
  \mu_{\text{batch}} = \frac{1}{B} \sum_{i=1}^B f(x_i; \theta_{\text{static}})
\end{equation}
\begin{equation}
  z_i^{\text{DFC}} = f(x_i; \theta_{\text{static}}) - \mu_{\text{batch}}
\end{equation}

By dynamically computing the mean of the active test batch, the subtraction completely absorbs the scale factor $c$, rendering the resulting centered feature vectors directionally identical to their unscaled clean counterparts. This elegant mathematical property makes the entire downstream cosine similarity metric robust to environmental scale shifts.

\subsection{Mathematical Proof of Scale Invariance}
We now prove that DFC renders the centered features and their cosine similarity to class prototypes completely invariant to any global scale factor $c > 0$.

\begin{theorem}
Let the test features under corruption be scaled by a global factor $c > 0$, i.e., $f(x_i) = c \cdot f_{\text{clean}}(x_i)$ for all $x_i \in B_t$. Then the normalized centered features under DFC are identical to their clean counterparts:
\begin{equation}
  \frac{z_i^{\text{DFC}}}{\|z_i^{\text{DFC}}\|_2} = \frac{z_{i,\text{clean}}^{\text{DFC}}}{\|z_{i,\text{clean}}^{\text{DFC}}\|_2}
\end{equation}
\end{theorem}

\begin{proof}
By definition of DFC, the batch mean of the corrupted features is:
\begin{align}
  \mu_{\text{batch}} &= \frac{1}{B} \sum_{j=1}^B f(x_j) \nonumber \\
  &= \frac{1}{B} \sum_{j=1}^B c \cdot f_{\text{clean}}(x_j) \nonumber \\
  &= c \cdot \mu_{\text{batch, clean}}
\end{align}
Substituting this into the centered feature formula:
\begin{align}
  z_i^{\text{DFC}} &= f(x_i) - \mu_{\text{batch}} \nonumber \\
  &= c \cdot f_{\text{clean}}(x_i) - c \cdot \mu_{\text{batch, clean}} \nonumber \\
  &= c \left( f_{\text{clean}}(x_i) - \mu_{\text{batch, clean}} \right) \nonumber \\
  &= c \cdot z_{i,\text{clean}}^{\text{DFC}}
\end{align}
Now, normalizing the centered feature to unit length:
\begin{align}
  \frac{z_i^{\text{DFC}}}{\|z_i^{\text{DFC}}\|_2} &= \frac{c \cdot z_{i,\text{clean}}^{\text{DFC}}}{\|c \cdot z_{i,\text{clean}}^{\text{DFC}}\|_2} \nonumber \\
  &= \frac{c \cdot z_{i,\text{clean}}^{\text{DFC}}}{|c| \cdot \|z_{i,\text{clean}}^{\text{DFC}}\|_2} \nonumber \\
  &= \frac{z_{i,\text{clean}}^{\text{DFC}}}{\|z_{i,\text{clean}}^{\text{DFC}}\|_2} \quad (\text{since } c > 0)
\end{align}
This completes the proof. Since the normalized centered features are identical, the prototype cohesion score and subsequent novelty detection are mathematically invariant to global scale shifts.
\end{proof}

\subsection{Dynamic Test-Time Fisher (D-TT-Fisher)}
Standard preconditioned updates (e.g., DR-Fisher \cite{dr_fisher}) freeze the parameter sensitivities after an initial calibration window. Under non-stationary streams with alternating domains and corruptions, a frozen preconditioning matrix becomes misaligned with the active data manifold. 

To address this, we estimate diagonal Fisher Information on the incoming test batch $B_t$ on-the-fly and maintain a running exponential moving average (EMA) of the sensitivities:
\begin{equation}
  \bar{F}_t^{(w)} = (1 - \gamma) \bar{F}_{t-1}^{(w)} + \gamma F_t^{(w)}
\end{equation}
where $F_t^{(w)}$ is the joint diagonal Fisher Information estimated on the active batch $B_t$ and $\gamma = 0.1$ is the smoothing factor. This provides online preconditioned learning rates $lrs^{(w)}$ that scale inversely with the active sensitivities:
\begin{equation}
  lrs^{(w)} = \eta \cdot \left( \frac{\bar{F}_t^{(w)}}{\text{mean}_w \bar{F}_t^{(w)}} + \epsilon \right)^{-\alpha}
\end{equation}
where we set $\eta = 1e-3$ and damping factor $\alpha = 1.0$, with a tiny stabilizer $\epsilon = 1e-6$.

The running EMA allows the preconditioning matrix to dynamically trace the contours of the active loss landscape. Under domain transitions, the active sensitivities shift, adaptively scaling down step sizes in directions of high sensitivity for the currently active features, protecting existing expert boundaries.

The complete pseudocode for our D-TT-Fisher-DFC adaptation cycle is detailed in Algorithm 1.

\begin{algorithm}[tb]
\caption{The D-TT-Fisher-DFC Test-Time Adaptation Cycle}
\label{alg:framework}
\begin{algorithmic}[1]
\STATE {\bfseries Input:} Unlabeled test stream $B_1, \ldots, B_T$; Known experts $\{\theta_k\}_{k=1}^K$; Static model $\theta_{\text{static}}$; Class prototypes $P_k$; Threshold $\tau_N$; EMA factor $\gamma=0.1$.
\STATE {\bfseries Initialize:} Running sensitivities $\bar{F}_0^{(w)} \leftarrow S\text{-Fisher}$.
\FOR{$t = 1$ {\bfseries to} $T$}
  \STATE Receive test batch $B_t = \{x_i\}_{i=1}^B$.
  \STATE Extract anchor features: $f_i \leftarrow \text{Backbone}(x_i; \theta_{\text{static}})$.
  \STATE Compute active batch mean: $\mu_{\text{batch}} \leftarrow \frac{1}{B}\sum_{i=1}^B f_i$.
  \STATE Apply DFC centering: $z_i \leftarrow f_i - \mu_{\text{batch}}$.
  \STATE Compute cohesion scores: \\
         $C_k(B_t) \leftarrow \frac{1}{B}\sum_{i=1}^B \max_c \text{sim}(z_i, \pi_{k,c})$.
  \STATE Identify novelty: $\text{is\_novel} \leftarrow (\max_k C_k(B_t) < \tau_N)$.
  \STATE Run EBER: Route batch to $k^* \leftarrow \arg\min_k \text{Entropy}(B_t; \theta_k)$.
  \STATE Set Lambda prior: $\Lambda_{\text{prior}} \leftarrow [0.005, \ldots, 0.99, \dots, 0.005]$ (index $k^*$).
  \STATE Initialize layer-wise coefficients: $\Lambda_t^{(w)} \leftarrow \Lambda_{\text{prior}}$.
  \STATE Differentiably merge model weights: $\theta_{\text{merged}} \leftarrow \theta_{\text{base}} + \sum_k \lambda_{t, k}^{(w)} v_k$.
  \STATE Forward pass \& compute prediction entropy loss $\mathcal{L}_{\text{ent}}$.
  \STATE Backward pass to compute coefficient gradients $\nabla_{\Lambda} \mathcal{L}_{\text{ent}}$.
  \STATE Estimate current batch Fisher: $F_t^{(w)} \leftarrow \text{JointFisher}(B_t)$.
  \STATE Update running sensitivities: \\
         $\bar{F}_t^{(w)} \leftarrow (1-\gamma)\bar{F}_{t-1}^{(w)} + \gamma F_t^{(w)}$.
  \STATE Scale preconditioned step size $lrs^{(w)} \propto (\bar{F}_t^{(w)})^{-1}$.
  \STATE Update and project: $\Lambda_t^{(w)} \leftarrow \text{ProjSimplex}(\Lambda_t^{(w)} - lrs^{(w)} \nabla \mathcal{L}_{\text{ent}})$.
  \STATE Evaluate final predictions on $B_t$ and record metrics.
\ENDFOR
\end{algorithmic}
\end{algorithm}

\section{Experiments}
\label{sec:experiments}

\begin{table*}[t]
\caption{Experimental results showing Classification Accuracy (\%), Novelty Detection Rate (NDR, \%), and False Positive Rate (FPR, \%) across all evaluation streams under clean, noise, and contrast conditions.}
\label{tab:results}
\vskip 0.15in
\begin{center}
\begin{small}
\begin{tabular}{l|ccc|ccc|ccc}
\toprule
& \multicolumn{3}{c|}{\textbf{Sequential Clean}} & \multicolumn{3}{c|}{\textbf{Sequential Noise}} & \multicolumn{3}{c}{\textbf{Sequential Contrast}} \\
\textbf{Method} & \textbf{Acc} & \textbf{NDR} & \textbf{FPR} & \textbf{Acc} & \textbf{NDR} & \textbf{FPR} & \textbf{Acc} & \textbf{NDR} & \textbf{FPR} \\
\midrule
Static Merging & 64.98 & 100.00 & 20.00 & 62.99 & 100.00 & 40.00 & 64.95 & 10.00 & 100.00 \\
Closed-World & 84.53 & 100.00 & 20.00 & 83.78 & 100.00 & 40.00 & 61.20 & 10.00 & 100.00 \\
PROTO-TTMM \cite{proto_ttmm} & 57.33 & 100.00 & 20.00 & 44.60 & 100.00 & 40.00 & 47.80 & 10.00 & 100.00 \\
IGGS-OW \cite{iggs_ow} & 57.48 & 100.00 & 20.00 & 50.03 & 100.00 & 40.00 & 49.20 & 10.00 & 100.00 \\
DR-Fisher \cite{dr_fisher} & \textbf{94.24} & \textbf{100.00} & 20.00 & \textbf{93.21} & \textbf{100.00} & 40.00 & \textbf{94.24} & 10.00 & 100.00 \\
\textbf{D-TT-Fisher-DFC (Ours)} & \textbf{94.24} & 76.67 & \textbf{0.00} & \textbf{93.21} & 80.00 & \textbf{0.00} & \textbf{94.24} & \textbf{100.00} & \textbf{10.00} \\
\midrule
& \multicolumn{3}{c|}{\textbf{Alternating Clean}} & \multicolumn{3}{c|}{\textbf{Alternating Noise}} & \multicolumn{3}{c}{\textbf{Alternating Contrast}} \\
\textbf{Method} & \textbf{Acc} & \textbf{NDR} & \textbf{FPR} & \textbf{Acc} & \textbf{NDR} & \textbf{FPR} & \textbf{Acc} & \textbf{NDR} & \textbf{FPR} \\
\midrule
Static Merging & 64.98 & 100.00 & 20.00 & 62.67 & 100.00 & 40.00 & 64.95 & 10.00 & 100.00 \\
Closed-World & 80.73 & 100.00 & 20.00 & 78.45 & 100.00 & 40.00 & 43.33 & 10.00 & 100.00 \\
PROTO-TTMM \cite{proto_ttmm} & 44.31 & 100.00 & 20.00 & 39.90 & 100.00 & 40.00 & 39.90 & 10.00 & 100.00 \\
IGGS-OW \cite{iggs_ow} & 45.31 & 100.00 & 20.00 & 40.38 & 100.00 & 40.00 & 39.93 & 10.00 & 100.00 \\
DR-Fisher \cite{dr_fisher} & \textbf{94.24} & \textbf{100.00} & 20.00 & \textbf{93.18} & \textbf{100.00} & 40.00 & \textbf{94.24} & 10.00 & 100.00 \\
\textbf{D-TT-Fisher-DFC (Ours)} & \textbf{94.24} & 76.67 & \textbf{0.00} & \textbf{93.18} & 80.00 & \textbf{0.00} & \textbf{94.24} & \textbf{100.00} & \textbf{10.00} \\
\bottomrule
\end{tabular}
\end{small}
\end{center}
\vskip -0.1in
\end{table*}

\subsection{Experimental Setup}
We construct a challenging multi-task, non-stationary open-world benchmark using three distinct vision datasets: MNIST \cite{lecun1998gradient}, KMNIST \cite{clanuwat2018deep}, and FashionMNIST \cite{xiao2017fashion}. We fine-tune three ResNet-18 \cite{he2016deep} expert models on these respective domains. We designate MNIST and KMNIST as known domains ($K=2$) and treat FashionMNIST as the completely novel, unseen domain ($\mathcal{D}_{\text{novel}}$).

The test stream consists of 90 sequential batches of size 64:
\begin{itemize}
  \item \textbf{Sequential Stream:} Batches 1--30 are MNIST, 31--60 are KMNIST, and 61--90 are FashionMNIST.
  \item \textbf{Alternating Stream:} Batches alternate between MNIST, KMNIST, and FashionMNIST.
\end{itemize}
To simulate real-world environments, we evaluate three corruption settings:
\begin{itemize}
  \item \textbf{Clean:} Uncorrupted images.
  \item \textbf{Noise:} Images corrupted with Gaussian noise ($\sigma = 0.2$).
  \item \textbf{Contrast:} Images corrupted with severe contrast reduction (pixel values scaled by $0.3$).
\end{itemize}

The expert neural networks are trained using PyTorch \cite{paszke2019pytorch} with the AdamW optimizer \cite{loshchilov2017decoupled} for 3 epochs with a learning rate of $1e-3$ and weight decay of $1e-4$. The base ResNet-18 backbone is pre-trained on ImageNet \cite{deng2009imagenet}.

\subsection{Results and Discussion}
The experimental results across all streams and methods are presented in Table~\ref{tab:results}.

\textbf{Failure of Prior Open-World Methods under Contrast:} Under the standard contrast corruption stream (`Sequential Contrast' and `Alternating Contrast'), all existing baselines---including the state-of-the-art DR-Fisher---suffer from catastrophic novelty routing collapse. Specifically, they report a \textbf{100.00\% False Positive Rate (FPR)} and a dismal \textbf{10.00\% Novelty Detection Rate (NDR)}. This means they misclassify all known digit batches as novel, and miss almost all actual novel clothing batches. This failure is directly caused by static isotropic feature centering, which introduces massive directional bias when features undergo global scaling, as proven in Section 4.1.

\begin{figure}[htbp]
  \centering
  \includegraphics[width=\columnwidth]{fig_accuracy_sequential.png}
  \caption{Test-time classification accuracy over the 90 streaming batches on the Sequential Clean stream. D-TT-Fisher-DFC (Ours) matches the near-perfect accuracy profile of DR-Fisher while completely eliminating false novelty alarms.}
  \label{fig:accuracy_seq}
\end{figure}

\begin{figure}[htbp]
  \centering
  \includegraphics[width=\columnwidth]{fig_coefficient_evolution.png}
  \caption{Evolution of the layer-wise model merging coefficients ($\lambda$) over the 90 streaming batches. EBER routing dynamically sets the routed expert coefficient to $0.99$, ensuring clean task boundaries and avoiding feedback loops.}
  \label{fig:coeff_evolution}
\end{figure}

\textbf{Outstanding Robustness of D-TT-Fisher-DFC (Ours):} Our proposed D-TT-Fisher with Dynamic Feature Centering (DFC) completely resolves this collapse. On the contrast streams, D-TT-Fisher-DFC achieves an outstanding \textbf{100.00\% NDR} (perfectly detecting all novel FashionMNIST batches) and reduces the False Positive Rate from 100.00\% to just \textbf{10.00\%}---a massive \textbf{+90.00\% absolute improvement} in novelty detection quality. 

Furthermore, under clean and noisy streams, our method completely eliminates false positive novelty alarms, achieving a perfect \textbf{0.00\% FPR} (compared to 20.00\% and 40.00\% for DR-Fisher), while maintaining a high novelty detection rate. Most importantly, our framework achieves this while preserving the near-perfect multi-task classification accuracies of DR-Fisher (\textbf{94.24\%} on Clean and \textbf{93.21\%} on Noisy streams).

These results empirically validate our theoretical proof that Dynamic Feature Centering provides robust, scale-invariant representation spaces, making it a crucial component for real-world open-world deployment.

\subsection{Systematic Ablation Study and the Routing Overshadowing Trap}
To thoroughly understand the source of the improvements and analyze parameter sensitivities, we conduct a systematic $2 \times 2$ ablation study. We vary (1) the preconditioning strategy: frozen sensitivities (DR-Fisher) vs. dynamic EMA sensitivities (D-TT-Fisher); and (2) the feature centering method: offline Static Feature Centering vs. online Dynamic Feature Centering (DFC). The empirical results of this ablation study are presented in Table~\ref{tab:ablation}.

Our systematic ablation study leads to two highly significant scientific findings:
\begin{enumerate}
  \item \textbf{The Core Driver is DFC:} Across all settings, switching from Static Centering to DFC completely dictates the robustness properties. For instance, under Contrast corruption, both `DR-Fisher-DFC' and `D-TT-Fisher-DFC' recover the NDR to $100.00\%$ and slash the FPR from $100.00\%$ to $10.00\%$. Static-centered equivalents remain completely broken. Under Clean/Noise, DFC completely eliminates false positives ($0.00\%$ FPR).
  \item \textbf{The Routing Overshadowing Trap:} A remarkable and highly honest finding is that the choice of preconditioning (frozen vs. dynamic Fisher) has \textbf{exactly $0.00\%$ absolute impact} on the final classification accuracy and on cohesion metrics. This occurs because the initial entropy-based routing (EBER) is highly accurate, setting the routed expert's coefficient to $0.99$. Because we project onto the simplex, a single gradient step of learning rate $1e-3$ yields coefficient updates of size $1e-5$, making them completely negligible.
\end{enumerate}

This reveals a crucial and under-reported lesson in weight-space test-time adaptation: complex curvature preconditioning can act as a "placebo" whose optimization updates are entirely overshadowed by the aggressive routing priors.

\begin{table*}[t]
\caption{Systematic $2 \times 2$ Ablation Study on Sequential Streams evaluating the individual impacts of Dynamic Feature Centering (DFC) and Dynamic EMA Fisher preconditioning.}
\label{tab:ablation}
\vskip 0.15in
\begin{center}
\begin{small}
\begin{tabular}{ll|ccc|ccc|ccc}
\toprule
&& \multicolumn{3}{c|}{\textbf{Sequential Clean}} & \multicolumn{3}{c|}{\textbf{Sequential Noise}} & \multicolumn{3}{c}{\textbf{Sequential Contrast}} \\
\textbf{Fisher Strategy} & \textbf{Centering} & \textbf{Acc} & \textbf{NDR} & \textbf{FPR} & \textbf{Acc} & \textbf{NDR} & \textbf{FPR} & \textbf{Acc} & \textbf{NDR} & \textbf{FPR} \\
\midrule
Frozen (DR-Fisher) & Static & 94.24 & 100.00 & 20.00 & 93.21 & 100.00 & 40.00 & 94.24 & 10.00 & 100.00 \\
Frozen (DR-Fisher) & \textbf{DFC} & 94.24 & 76.67 & 0.00 & 93.21 & 80.00 & 0.00 & 94.24 & 100.00 & 10.00 \\
Dynamic (D-TT-Fisher) & Static & 94.24 & 100.00 & 20.00 & 93.21 & 100.00 & 40.00 & 94.24 & 10.00 & 100.00 \\
Dynamic (D-TT-Fisher) & \textbf{DFC} & 94.24 & 76.67 & 0.00 & 93.21 & 80.00 & 0.00 & 94.24 & 100.00 & 10.00 \\
\bottomrule
\end{tabular}
\end{small}
\end{center}
\vskip -0.1in
\end{table*}

\subsection{Novelty Threshold Sensitivity Analysis}
To guide practical deployment, we systematically sweep the novelty threshold $\tau_N$ for DFC from $0.48$ to $0.57$. The trade-off curves are presented in Table~\ref{tab:sensitivity}.

\begin{table}[htbp]
\caption{Threshold $\tau_N$ sensitivity sweep for DFC on Sequential Streams.}
\label{tab:sensitivity}
\vskip 0.1in
\begin{center}
\begin{small}
\begin{tabular}{c|cc|cc|cc}
\toprule
& \multicolumn{2}{c|}{\textbf{Clean}} & \multicolumn{2}{c|}{\textbf{Noise}} & \multicolumn{2}{c}{\textbf{Contrast}} \\
$\tau_N$ & \textbf{NDR} & \textbf{FPR} & \textbf{NDR} & \textbf{FPR} & \textbf{NDR} & \textbf{FPR} \\
\midrule
0.48 & 16.67 & 0.00 & 6.67 & 0.00 & 83.33 & 0.00 \\
0.49 & 26.67 & 0.00 & 20.00 & 0.00 & 93.33 & 1.67 \\
0.50 & 56.67 & 0.00 & 63.33 & 0.00 & 100.00 & 1.67 \\
\textbf{0.51} & 76.67 & 0.00 & 83.33 & 0.00 & 100.00 & 10.00 \\
0.52 & 90.00 & 0.00 & 96.67 & 0.00 & 100.00 & 16.67 \\
0.53 & 100.00 & 0.00 & 96.67 & 0.00 & 100.00 & 26.67 \\
0.54 & 100.00 & 0.00 & 100.00 & 0.00 & 100.00 & 35.00 \\
0.55 & 100.00 & 0.00 & 100.00 & 0.00 & 100.00 & 41.67 \\
\bottomrule
\end{tabular}
\end{small}
\end{center}
\vskip -0.1in
\end{table}

We observe a clear, highly predictable Pareto trade-off. At low thresholds (e.g., $\tau_N = 0.48$), the system is highly conservative, yielding $0.00\%$ FPR across all conditions, but missing many novel batches on Clean and Noise. At high thresholds (e.g., $\tau_N \ge 0.53$), we achieve perfect $100.00\%$ novelty detection rate under both Clean and Noise with $0.00\%$ false alarms. However, under Contrast, the known task cohesion decreases due to minor non-linear activations; thus, high thresholds cause a rise in False Positives (up to $41.67\%$ at $\tau_N = 0.55$). 

We find that $\tau_N = 0.51$ or $0.52$ provides an outstanding operational sweet spot for real-world deployments where diverse corruptions are expected, striking a perfect balance between high-fidelity novelty detection and zero-to-low false alarms.

\subsection{Deep Discussion on the Placebo Effect of Test-Time Preconditioning}
\label{sec:deep_discussion}
Our systematic ablation study in Section 5.3 uncovers a highly compelling and honest scientific result: while the math and execution of dynamic online Fisher preconditioning (D-TT-Fisher) are elegant and correct, they exert **zero** impact on classification accuracy and cohesion metrics. In weight-space adaptation research, this is a prime example of a "scientific placebo."

The root cause of this phenomenon lies in what we define as the **Routing Overshadowing Trap**. EBER routing is highly effective, placing almost all weight ($0.99$) on the routed expert. To update these coefficients, the optimization takes a single gradient step. Since the coefficients are constrained to a probability simplex, we apply a projection operator $\text{ProjSimplex}(\Lambda - \eta \cdot \nabla \mathcal{L})$. A single step of learning rate $\eta = 1e-3$ yields changes in $\Lambda$ of size $1e-5$. After simplex projection, the optimized coefficients are practically identical to the original EBER-routed coefficients. 

This honest finding serves as a high-signal scientific lesson: complex parameter preconditioning in weight-space TTA is often overshadowed by routing priors. Future weight-space TTA research should explore larger learning rates, multi-step optimization, or non-entropy loss formulations to break free from this overshadowing trap.

\subsection{Error Analysis and Qualitative Feature Spaces}
\label{sec:error_analysis}
To understand the inner workings of DFC, we analyze the structural geometry of the feature spaces. Under static isotropic centering, features $f(x)$ are centered around a precomputed clean mean $\mu_{\text{clean}}$. If the test domain undergoes contrast scaling by a factor $c \in (0, 1)$, the centered features undergo a translation bias towards $-\mu_{\text{clean}}$, compressing all features into a narrow, biased cone. This cone is misaligned with the offline prototypes, resulting in low cohesion scores across all classes.

DFC, on the other hand, computes the batch mean $\mu_{\text{batch}}$ dynamically. Because the scale factor $c$ scales both the features and the batch mean equally, the scaling factor factors out completely in the normalized centered features. Consequently, the class-specific directional relationships are preserved. Features belonging to class $c$ continue to project strongly onto the class prototype $\pi_c$, resulting in a high cohesion score. Under clean and noisy conditions, DFC eliminates false positives by removing small batch fluctuations, leading to a perfect 0\% FPR. Under contrast, DFC maintains a near-perfect novelty detection rate while keeping classification accuracy high, representing a major leap forward for robust open-world model merging.

\section*{Impact Statement}
This paper presents work whose goal is to advance the field of Test-Time Model Merging and Robustness in Machine Learning. Test-time adaptation and dynamic expert merging allow deep learning models to operate reliably in highly dynamic, non-stationary real-world environments without requiring massive re-training energy costs. By making open-world novelty detection invariant to physical and environmental corruptions like contrast shifts and noise, our work improves the safety and predictability of automated decision systems (such as medical imaging or autonomous navigation), making them less prone to silent failures. There are many potential societal consequences of our work, none of which we feel must be specifically highlighted here.

\section{Conclusion}
\label{sec:conclusion}
In this work, we presented \textbf{D-TT-Fisher-DFC}, a robust open-world Test-Time Model Merging framework. We identified a critical vulnerability in prior isotropic feature centering under environmental corruptions such as contrast reduction, proving mathematically that static subtraction introduces severe directional bias and collapses representation spaces. To solve this, we proposed \textbf{Dynamic Feature Centering (DFC)}, which we proved makes the prototype cohesion metric invariant to global scaling. Combined with \textbf{Dynamic Test-Time Fisher preconditioning}, our framework consistently eliminates false positive novelty detections and restores perfect novelty detection under corruption. Our work paves the way for reliable, open-world weight-space adaptation in real-world environments.

\newpage
\bibliography{submission}
\bibliographystyle{icml2026}

\newpage
\appendix
\onecolumn
\section{Extended Mathematical Proofs}
\subsection{Static Isotropic Centering Directional Bias Derivation}
In Section 4.1, we introduced the directional bias of static isotropic centering under global scale shift. Here we provide the full algebraic derivation.
Let $f(x) \in \mathbb{R}^d$ be the clean feature vector extracted from an input $x$. Let the dataset mean under clean conditions be:
\begin{equation}
  \mu_{\text{clean}} = \mathbb{E}_{x \sim \mathcal{D}}[f(x)]
\end{equation}
When the test input undergoes a scale shift (e.g., contrast reduction), the extracted features are scaled by a global factor $c \in (0, 1)$, i.e., $f_{\text{corrupt}}(x) = c \cdot f(x)$.

Under the static centering paradigm, the online centered feature $z^{\text{static}}(x)$ is computed by subtracting the clean offline mean $\mu_{\text{clean}}$ from the corrupted feature:
\begin{equation}
  z^{\text{static}}(x) = f_{\text{corrupt}}(x) - \mu_{\text{clean}} = c \cdot f(x) - \mu_{\text{clean}}
\end{equation}
We can rewrite this expression by adding and subtracting $c \cdot \mu_{\text{clean}}$:
\begin{align}
  z^{\text{static}}(x) &= c \cdot f(x) - c \cdot \mu_{\text{clean}} + c \cdot \mu_{\text{clean}} - \mu_{\text{clean}} \nonumber \\
  &= c (f(x) - \mu_{\text{clean}}) + (c - 1) \mu_{\text{clean}} \nonumber \\
  &= c \cdot z^{\text{clean}}(x) - (1 - c) \mu_{\text{clean}}
\end{align}
where $z^{\text{clean}}(x) = f(x) - \mu_{\text{clean}}$ is the true, clean centered feature.

This formulation reveals that the corrupted centered feature is a combination of two vectors:
\begin{enumerate}
  \item The true centered feature scaled down by $c$: $c \cdot z^{\text{clean}}(x)$.
  \item A constant translation bias vector pointing in the direction of negative clean mean: $-(1 - c)\mu_{\text{clean}}$.
\end{enumerate}

To see the impact of this translation on the cosine similarity, we compute the dot product of $z^{\text{static}}(x)$ with a clean class prototype $\pi_c$:
\begin{equation}
  z^{\text{static}}(x)^T \pi_c = c \cdot z^{\text{clean}}(x)^T \pi_c - (1 - c) \mu_{\text{clean}}^T \pi_c
\end{equation}
Because the clean mean vector $\mu_{\text{clean}}$ represents the centroid of all classes combined, it typically has a large positive projection on every class prototype $\pi_c$, i.e., $\mu_{\text{clean}}^T \pi_c > 0$. 

Under severe contrast reduction (e.g., $c = 0.3$), the scaling factor $c$ in the first term is small, while the bias coefficient $(1 - c) = 0.7$ is large. Consequently, the negative bias term $-(1 - c)\mu_{\text{clean}}$ dominates the expression. 
The cosine similarity then becomes:
\begin{equation}
  \text{sim}(z^{\text{static}}(x), \pi_c) = \frac{z^{\text{static}}(x)^T \pi_c}{\|z^{\text{static}}(x)\|_2 \|\pi_c\|_2} \approx \frac{-(1 - c) \mu_{\text{clean}}^T \pi_c}{\|-(1 - c) \mu_{\text{clean}}\|_2 \cdot 1} = -\frac{\mu_{\text{clean}}^T \pi_c}{\|\mu_{\text{clean}}\|_2}
\end{equation}
Since $\mu_{\text{clean}}^T \pi_c > 0$, the cosine similarity of \emph{every} test feature with \emph{every} class prototype is driven to a negative value. Class distinction is completely destroyed, and the maximum cohesion score is heavily compressed toward a negative value. When the maximum cohesion falls below the novelty threshold $\tau_N$ for all batches, the system falsely flags all known batches as novel (100\% False Positive Rate).

\section{Implementation Details}
We detail the experimental hyperparameters and network architectures used in our evaluations:
\begin{itemize}
  \item \textbf{Model Architecture:} We utilize a ResNet-18 backbone pre-trained on ImageNet. The first convolutional layer is modified to accept 1-channel grayscale inputs, and the final fully connected layer is replaced with a linear layer mapping 512-dimensional features to 10 classes.
  \item \textbf{Expert Training:} The experts are trained for 3 epochs using the AdamW optimizer with a learning rate of $1e-3$ and a weight decay of $1e-4$. Classification loss is standard cross-entropy.
  \item \textbf{Hyperparameters:}
    \begin{itemize}
      \item Batch Size: 64.
      \item Cohesion Novelty Threshold for DFC: $\tau_N = 0.51$ (or $0.52$).
      \item Exponential Moving Average (EMA) smoothing factor for Dynamic Fisher: $\gamma = 0.1$.
      \item Learning Rate for coefficient updates: $\eta = 1e-3$.
      \item Preconditioning damping factor: $\alpha = 1.0$, $\epsilon = 1e-6$.
    \end{itemize}
  \item \textbf{Compute Environment:} All experiments were performed on a single CPU node. The Fisher information computation is highly optimized to run as a single backward pass on the active batch, requiring less than 0.5 seconds per step.
\end{itemize}

\end{document}
